% Certification and frontier chapter.  Every correctness assertion below is
% accompanied by a proof; conjectures are confined to the final frontier
% section and include explicit proof obligations.

% LOCAL-NON-FOURIER-HAT: \CertifiedEnclosureCenter denotes a certified numerical enclosure center approximating the exact quantity supplied as the argument.
\newcommand{\CertifiedEnclosureCenter}[1]{\widehat{#1}}
% LOCAL-NON-FOURIER-HAT: \FormalGevreySeriesSymbol denotes a formal Gevrey series before Borel transformation, not a Fourier transform.
\newcommand{\FormalGevreySeriesSymbol}[1]{\widehat{#1}}

\part{Certification and research frontiers}

\section{Branch-aware numerical certification}
\label{sec:certification-frontiers}

An inverse approximation is not a certificate.  A certificate must identify
the forward map, the active parameter, the target, the domain containing the
selected branch, and the arithmetic enclosure used to control every discarded
tail and rounding error.  In the complex setting it must additionally record
the continuation path and every logarithm or fractional-power determination.
Changing any one of these data can change the inverse being computed.

The basic workflow is therefore:
\begin{enumerate}
\item choose a coordinate in which the selected sheet is regular;
\item evaluate the forward map and its derivatives with outward-rounded
      enclosures, including explicit infinite tails;
\item obtain a candidate from an asymptotic expansion, continuation, or a
      floating-point solve;
\item prove existence and uniqueness in a real interval or complex disk;
\item if all sheets are required, certify the remaining zero count by the
      argument principle.
\end{enumerate}
The candidate-generation step may be heuristic.  The last two steps are not.

\subsection{Certified tails for infinite products}
\label{sec:certified-product-tails}

The next estimate controls the value and the first two argument derivatives
with one truncation index.

\begin{theorem}[Complex \(q\)-Pochhammer tail]
\label{thm:qpoch-tail-certificate}
Let \(a,q\in\ComplexNumbers\), put \(\rho=\lvert q\rvert<1\), and choose \(N\in\NaturalNumbers\) so that
\[
  \eta:=\lvert a\rvert\rho^N<1.
\]
The tail
\[
  T_N(a,q):=\prod_{j=N}^{\infty}(1-aq^j)
\]
is nonzero and has the logarithm
\[
  L_N(a,q):=-\sum_{j=N}^{\infty}\sum_{k=1}^{\infty}
                 \frac{(aq^j)^k}{k}.
\]
With
\[
  B_N:=\frac{\lvert a\rvert\rho^N}
              {(1-\eta)(1-\rho)},
\]
one has
\begin{align}
 \lvert L_N(a,q)\rvert&\le B_N,
 &\lvert T_N(a,q)-1\rvert&\le \EulerE^{B_N}-1,
 \label{eq:qpoch-tail-value}\\
 \lvert\partial_aL_N(a,q)\rvert
   &\le\frac{\rho^N}{(1-\eta)(1-\rho)},
 &\lvert\partial_a^2L_N(a,q)\rvert
   &\le\frac{\rho^{2N}}{(1-\eta)^2(1-\rho^2)}.
 \label{eq:qpoch-tail-derivatives}
\end{align}
More explicitly, suppose interval arithmetic certifies scalar bounds
\[
 \rho\le\bar\rho<1,\qquad \eta\le\bar\eta<1.
\]
Then the value bounds hold with
\[
 \bar B_N:=\frac{\bar\eta}
 {(1-\bar\eta)(1-\bar\rho)}
\]
in place of \(B_N\), and the two derivative bounds hold after replacing
their right-hand sides respectively by
\[
 \frac{\bar\rho^N}{(1-\bar\eta)(1-\bar\rho)},
 \qquad
 \frac{\bar\rho^{2N}}{(1-\bar\eta)^2(1-\bar\rho^2)}.
\]
For interval balls \(\mathcal A\) and \(\mathcal Q\) of possible \(a\) and
\(q\), it suffices to certify
\[
 \bar\rho\ge\sup_{q\in\mathcal Q}\lvert q\rvert,\qquad
 \bar\eta\ge
   \left(\sup_{a\in\mathcal A}\lvert a\rvert\right)\bar\rho^N,
\]
with both upper endpoints still strictly below one.
\end{theorem}

\begin{proof}
For \(\lvert z\rvert<1\), the power-series logarithm gives
\[
 \left\lvert-\log(1-z)\right\rvert
 \le\sum_{k\ge1}\lvert z\rvert^k
 =\frac{\lvert z\rvert}{1-\lvert z\rvert}.
\]
For \(j\ge N\), \(\lvert aq^j\rvert\le\eta\).  Hence
\[
 \lvert L_N(a,q)\rvert
 \le \sum_{j\ge N}\frac{\lvert a\rvert\rho^j}{1-\eta}
 =B_N.
\]
The majorant is summable, so the logarithmic series converges absolutely and
the product converges to \(T_N=\EulerE^{L_N}\).  No tail factor vanishes because
\(\lvert aq^j\rvert<1\).  The elementary estimate
\(\lvert\EulerE^z-1\rvert\le\EulerE^{\lvert z\rvert}-1\) proves the second bound in
\eqref{eq:qpoch-tail-value}.

On a slightly smaller neighbourhood of the given \(a\), the differentiated
series have the summable majorants used below; termwise differentiation is
therefore legitimate.  Explicitly,
\[
 \partial_aL_N=-\sum_{j\ge N}\frac{q^j}{1-aq^j},
 \qquad
 \partial_a^2L_N=-\sum_{j\ge N}\frac{q^{2j}}{(1-aq^j)^2}.
\]
Using \(\lvert1-aq^j\rvert\ge1-\eta\) and summing the two geometric series
gives \eqref{eq:qpoch-tail-derivatives}.  Each factor in the three displayed
majorants is monotone in \(\rho\) and \(\eta\) on \([0,1)\), so substituting
\(\bar\rho,\bar\eta\) proves the interval versions.
\end{proof}

Near a zero of a finite prefix, that factor should be removed symbolically.
The theorem then controls the nonvanishing tail without hiding the local zero
inside a logarithmic sum.

The Rvachev Fourier product has an equally simple real-axis certificate.  We
use \(\SincRad x=\sin(x)/x\) for \(x\ne0\) and \(\SincRad(0)=1\).

\begin{theorem}[Geometric sinc-product tail]
\label{thm:sinc-tail-certificate}
Let \(0<q<1\), \(t\in\RealNumbers\), and \(w_j=(1-q)q^j\).  Then
\[
 S_N(t):=\prod_{j=N}^{\infty}\SincRad(w_jt)
\]
converges, and
\begin{equation}
 \lvert S_N(t)-1\rvert
 \le \exp\!\left(
   \frac{t^2(1-q)^2q^{2N}}{6(1-q^2)}
 \right)-1.
 \label{eq:sinc-tail-certificate}
\end{equation}
\end{theorem}

\begin{proof}
Taylor's formula with integral remainder gives, for real \(x\),
\[
 \lvert\sin x-x\rvert
 \le\int_0^{\lvert x\rvert}(\lvert x\rvert-u)u\,\Differential u
 =\frac{\lvert x\rvert^3}{6}.
\]
Consequently \(\lvert\SincRad x-1\rvert\le x^2/6\), with the value at zero
understood by continuity.  Put \(u_j=\SincRad(w_jt)-1\).  For a finite product,
\[
 \left\lvert\prod_{j=N}^{M}(1+u_j)-1\right\rvert
 \le\prod_{j=N}^{M}(1+\lvert u_j\rvert)-1
 \le\exp\!\left(\sum_{j=N}^{M}\lvert u_j\rvert\right)-1.
\]
Moreover,
\[
 \sum_{j\ge N}\lvert u_j\rvert
 \le\frac{t^2}{6}\sum_{j\ge N}w_j^2
 =\frac{t^2(1-q)^2q^{2N}}{6(1-q^2)}.
\]
The summability proves convergence of the product, and passage to the limit in
the finite-product inequality proves \eqref{eq:sinc-tail-certificate}.
\end{proof}

\subsection{Real branches: brackets, residuals, and interval Newton}
\label{sec:real-branch-certification}

On a real branch, monotonicity is the cleanest existence and uniqueness
certificate.  Newton acceleration may be used inside the bracket, but it is
the bracket that carries the proof.

\begin{theorem}[Monotone branch certificate]
\label{thm:monotone-branch-certificate}
Let \(f:[\alpha,\beta]\to\RealNumbers\) be continuous and strictly monotone, and let
\(y\) lie between \(f(\alpha)\) and \(f(\beta)\).  Then there is a unique
\(x_*\in[\alpha,\beta]\) with \(f(x_*)=y\).  Sign-preserving bisection produces
nested brackets \(X_r\) containing \(x_*\) with
\[
 \operatorname{length}(X_r)=2^{-r}(\beta-\alpha).
\]
If \(f\) is differentiable and
\(\lvert f'(x)\rvert\ge m>0\) throughout the original bracket, then every
\(x_0\in[\alpha,\beta]\) satisfies
\begin{equation}
 \lvert x_0-x_*\rvert\le
 \frac{\lvert f(x_0)-y\rvert}{m}.
 \label{eq:real-residual-transfer}
\end{equation}
\end{theorem}

\begin{proof}
The intermediate value theorem gives existence, and strict monotonicity gives
uniqueness.  At each bisection step, one of the two half-intervals still has
endpoint values bracketing \(y\); choosing it preserves \(x_*\) and halves the
length.  Finally, the mean-value theorem gives
\[
 f(x_0)-f(x_*)=f'(\xi)(x_0-x_*)
\]
for some \(\xi\) between \(x_0\) and \(x_*\).  Taking absolute values and using
the derivative lower bound proves \eqref{eq:real-residual-transfer}.
\end{proof}

The residual estimate assumes that the root already exists.  Interval Newton
can certify existence as well, provided its derivative enclosure is separated
from zero.

\begin{theorem}[One-dimensional interval Newton]
\label{thm:interval-newton-one-dimensional}
Let \(X=[\alpha,\beta]\), let \(f\) be continuously differentiable on \(X\),
and let \(D\) be a nonempty closed bounded real interval containing \(f'(X)\)
with \(0\notin D\).
For \(x_0\in X\), define
\[
 N(X):=x_0-\frac{f(x_0)-y}{D}
 =\left\{x_0-\frac{f(x_0)-y}{d}:d\in D\right\}.
\]
If \(N(X)\subseteq X\), then \(f(x)=y\) has exactly one solution in \(X\).
If \(N(X)\subseteq\operatorname{int}X\), the solution lies in the interior of
\(X\).
\end{theorem}

\begin{proof}
First suppose \(D=[m,M]\) with \(0<m\le M\).  Then \(f'\ge m\), so \(f\) is
strictly increasing.  Put \(r=f(x_0)-y\).  If \(r=0\), \(x_0\) is the unique
root.  If \(r>0\), the left endpoint of \(N(X)\) is \(x_0-r/m\).  Inclusion in
\(X\) gives \(x_0-\alpha\ge r/m\).  The mean-value theorem therefore yields
\[
 f(x_0)-f(\alpha)\ge m(x_0-\alpha)\ge r,
\]
so \(f(\alpha)\le y<f(x_0)\), and the intermediate value theorem gives a
root.  If \(r<0\), the right endpoint of \(N(X)\) is \(x_0-r/m\); inclusion
gives \(\beta-x_0\ge-r/m\), whence
\(f(\beta)-f(x_0)\ge-r\).  Thus \(f(x_0)<y\le f(\beta)\), and again a root
exists.  Strict increase makes it unique.

If \(D\) is negative, apply the preceding argument to \(-f\) and target
\(-y\); its interval Newton set is the same.  The stronger inclusion
\(N(X)\subseteq\operatorname{int}X\) makes the relevant endpoint inequality
strict, and therefore excludes an endpoint root when \(r\ne0\).  When
\(r=0\), it says directly that the root \(x_0\in N(X)\) is interior.
\end{proof}

Directed rounding is essential here.  A derivative sampled at many points is
not an interval extension of \(f'(X)\), and an ordinary floating-point
inclusion is not a mathematical inclusion.

\subsection{Complex disks: contraction, Rouch\'e, and total zero counts}
\label{sec:complex-branch-certification}

The complex analogue of interval Newton is particularly transparent in one
variable.  It certifies a selected sheet and gives a constructive refinement
map.

\begin{theorem}[Complex Krawczyk disk]
\label{thm:complex-krawczyk-disk}
Let \(f\) be holomorphic on a neighbourhood of the closed disk
\(D=\{z:\lvert z-z_0\rvert\le r\}\), where \(r>0\), and let
\(A\in\ComplexNumbers\setminus\{0\}\).  If
\[
 \kappa:=\sup_{z\in D}\lvert1-Af'(z)\rvert<1,
 \qquad
 \lvert Af(z_0)\rvert\le(1-\kappa)r,
\]
then \(f\) has exactly one zero in \(D\).  If the second inequality is strict,
the zero lies in the open disk.
\end{theorem}

\begin{proof}
Set \(T(z)=z-Af(z)\).  The disk is convex, and integration along the line
segment from \(w\) to \(z\) gives
\[
 \lvert T(z)-T(w)\rvert
 \le\kappa\lvert z-w\rvert.
\]
Also, for \(z\in D\),
\[
 \lvert T(z)-z_0\rvert
 \le\lvert T(z)-T(z_0)\rvert+\lvert Af(z_0)\rvert
 \le\kappa r+(1-\kappa)r=r.
\]
Thus \(T\) is a contraction of the complete metric space \(D\) into itself.
The contraction theorem supplies a unique fixed point \(z_*\).  Since
\(A\ne0\), \(T(z_*)=z_*\) is equivalent to \(f(z_*)=0\).  Every zero is a fixed
point, so there are no others.  Under the strict second inequality the
displayed estimate is strict on the boundary, hence \(z_*\) is interior.
\end{proof}

When a second-derivative bound is sharper than an enclosure of the Newton map,
the following Rouch\'e certificate is convenient.

\begin{theorem}[One-zero Rouch\'e disk]
\label{thm:rouche-one-zero-disk}
Let \(f\) be holomorphic on a neighbourhood of
\(D=\{z:\lvert z-z_0\rvert\le r\}\), where \(r>0\), and suppose \(M\ge0\)
and
\[
 \sup_{z\in D}\lvert f''(z)\rvert\le M,
 \qquad
 \lvert f(z_0)\rvert+\frac12Mr^2<\lvert f'(z_0)\rvert r.
\]
Then \(f\) has exactly one zero, counted with multiplicity, in the open disk.
\end{theorem}

\begin{proof}
For \(z\in D\), Taylor's formula along the segment from \(z_0\) to \(z\) is
\[
 f(z)=f(z_0)+f'(z_0)(z-z_0)
 +(z-z_0)^2\int_0^1(1-s)f''(z_0+s(z-z_0))\,\Differential s.
\]
The last term has modulus at most \(Mr^2/2\) on the boundary.  Therefore, on
\(\lvert z-z_0\rvert=r\), the difference between \(f(z)\) and the linear
function \(f'(z_0)(z-z_0)\) has modulus strictly smaller than
\(\lvert f'(z_0)(z-z_0)\rvert\).  Rouch\'e's theorem says that the two
functions have the same number of zeros in the disk.  The linear function has
one, counted with multiplicity.
\end{proof}

Disjoint one-zero disks do not by themselves prove that every inverse sheet
has been found.  A validated argument-principle count closes that gap.

\begin{theorem}[Validated argument-principle count]
\label{thm:validated-argument-principle}
Let \(\Omega\) be a bounded Jordan domain, let \(f\) be holomorphic on a
neighbourhood of \(\overline\Omega\), and suppose interval evaluation proves
\(0\notin f(\partial\Omega)\).  Write the positively oriented boundary as a
finite concatenation of \(C^3\) panels
\(\gamma_k:[0,1]\to\partial\Omega\), and put
\[
 g_k(s):=\frac{f'(\gamma_k(s))\gamma_k'(s)}{f(\gamma_k(s))}.
\]
Suppose validated bounds
\(M_k\ge\sup_{0\le s\le1}\lvert g_k''(s)\rvert\) are available.  Let
\[
 Q:=\sum_k\frac{g_k(0)+g_k(1)}2,
 \qquad
 E:=\frac1{12}\sum_kM_k.
\]
Suppose outward-rounded arithmetic provides
\(\CertifiedEnclosureCenter{Q}\in\ComplexNumbers\) and \(\varepsilon_Q\ge0\) with
\(\lvert Q-\CertifiedEnclosureCenter{Q}\rvert\le\varepsilon_Q\).  If the closed disk
\[
\left\{z:\left\lvert z-\frac{\CertifiedEnclosureCenter{Q}}{2\pi\ImaginaryUnit}\right\rvert
       \le\frac{E+\varepsilon_Q}{2\pi}\right\}
\]
contains exactly one integer \(n\), then \(f\) has exactly \(n\) zeros in
\(\Omega\), counted with multiplicity.
\end{theorem}

\begin{proof}
For a twice continuously differentiable complex-valued function \(g\) on
\([0,1]\), two integrations by parts give the trapezoidal remainder formula
\[
 \int_0^1g(s)\,\Differential s-\frac{g(0)+g(1)}2
 =-\frac12\int_0^1s(1-s)g''(s)\,\Differential s.
\]
Since \(\int_0^1s(1-s)\,\Differential s=1/6\), its modulus is at most
\(\sup\lvert g''\rvert/12\).  Summing over the panels gives
\[
 \left\lvert\int_{\partial\Omega}\frac{f'(z)}{f(z)}\,\Differential z-Q\right\rvert
 \le E.
\]
The boundary nonvanishing hypothesis makes the argument principle applicable,
so
\[
 \frac1{2\pi\ImaginaryUnit}\int_{\partial\Omega}\frac{f'(z)}{f(z)}\,\Differential z
\]
is the integer equal to the number of zeros of \(f\) in \(\Omega\).  The
quadrature and rounding bounds, followed by the triangle inequality, place
this integer in the displayed disk.  If that disk contains only \(n\), the
zero count is \(n\).
\end{proof}

In practice, interval balls for \(f(\gamma_k([0,1]))\) certify boundary
nonvanishing, automatic differentiation bounds \(g_k''\), and panel
subdivision reduces \(E\).  Counting one zero in each selected disk and then
the same total number in a containing contour certifies both uniqueness and
completeness.  For a meromorphic forward map, the same integral counts zeros
minus poles; the poles and their multiplicities must be enclosed and added
back explicitly.

\subsection{Asymptotic truncation is part of the certificate}
\label{sec:asymptotic-truncation-discipline}

An \(O(t^{N+1})\) printed in a manuscript has no numerical content until its
constant, domain, and branch are known.  For a convergent Taylor expansion,
Cauchy's estimate supplies all three.

\begin{theorem}[Cauchy tail bound]
\label{thm:cauchy-truncation-tail}
Let \(H\) be holomorphic on a neighbourhood of
\(\{z:\lvert z\rvert\le R\}\), let
\(M\ge\max_{\lvert z\rvert=R}\lvert H(z)\rvert\), and write
\(H(z)=\sum_{j\ge0}c_jz^j\).  For \(N\in\NaturalNumbers\), \(0\le r<R\), and
\(\lvert z\rvert\le r\),
\[
 \left\lvert H(z)-\sum_{j=0}^{N}c_jz^j\right\rvert
 \le M\frac{(r/R)^{N+1}}{1-r/R}.
\]
\end{theorem}

\begin{proof}
Cauchy's coefficient formula gives \(\lvert c_j\rvert\le M/R^j\).  Hence
\[
 \left\lvert\sum_{j>N}c_jz^j\right\rvert
 \le M\sum_{j>N}(r/R)^j
 =M\frac{(r/R)^{N+1}}{1-r/R}.
\]
\end{proof}

For an inverse expansion, it is usually sharper to certify its forward
residual than to bound inverse coefficients separately.

\begin{theorem}[Branchwise inverse residual transfer]
\label{thm:inverse-residual-transfer}
For each parameter \(s\), let \(p(s)\) be the selected real solution of
\(F_s(p)=y_s\) in an interval \(I_s\), and suppose
\(p\mapsto F_s(p)\) is continuously differentiable there.  Let
\(p_N(s)\in I_s\) be any truncated inverse approximation.  If
\[
 \inf_{p\in I_s}\lvert\partial_pF_s(p)\rvert\ge m_s>0,
 \qquad
 \lvert F_s(p_N(s))-y_s\rvert\le E_N(s),
\]
then
\[
 \lvert p_N(s)-p(s)\rvert\le\frac{E_N(s)}{m_s}.
\]
\end{theorem}

\begin{proof}
Apply the mean-value theorem to \(p_N(s)\) and \(p(s)\) and use the derivative
lower bound.  The fact that both points lie on the certified interval is what
makes the bound branch-specific.
\end{proof}

At a fold, \(m_s\) tends to zero in the original coordinate.  One must first
pass to the square-root coordinate in which the sheet is regular, or use a
Rouch\'e or Krawczyk disk.  Reporting a large residual quotient in the singular
coordinate does not certify the branch.

Divergent asymptotic series require a different discipline.  The elementary
calculation below explains optimal truncation for a Gevrey-one envelope.

\begin{proposition}[Optimal index for a Gevrey-one remainder envelope]
\label{prop:gevrey-least-term}
Let \(A,C>0\) and \(t\ne0\).  For every \(N\in\NaturalNumbers\), define the candidate
remainder envelope
\[
 E_N(t)=C(A\lvert t\rvert)^{N+1}(N+1)!.
\]
Suppose the actual remainder after truncation at \(N\) has been proved to
have modulus at most \(E_N(t)\) for every \(N\).  Then
\[
 \frac{E_{N+1}(t)}{E_N(t)}=A\lvert t\rvert(N+2).
\]
Put \(x=(A\lvert t\rvert)^{-1}-2\).  The certified envelope decreases
exactly when \(N<x\), is unchanged when \(N=x\), and increases when \(N>x\).
If \(A\lvert t\rvert>1/2\), its unique minimizing index is \(N=0\).  If
\(A\lvert t\rvert=1/2\), the minimizing indices are \(0\) and \(1\).  If
\(A\lvert t\rvert<1/2\) and \(x\notin\NaturalNumbers\), the unique minimizing index is
\(\Ceiling{x}\); if \(x\in\NaturalNumbers\), the minimizing indices are \(x\) and
\(x+1\).
\end{proposition}

\begin{proof}
Cancel the common factors in \(E_{N+1}/E_N\).  The monotonicity and location of
the discrete minimum follow immediately from whether this ratio is below or
above one.
\end{proof}

This proposition is not permission to assume a Gevrey bound.  Such a bound
must first be proved uniformly in the stated sector.  A defensible asymptotic
certificate therefore records the coordinate and sheet, the allowed domain,
the remainder constants, the chosen truncation order, the outward-rounded
forward residual, and the final real or complex root enclosure.  Merely
observing that successive terms initially decrease is evidence, not proof.

\section{Two archived conjectures that are theorems}
\label{sec:upgraded-conjectures}

Two claims repeatedly presented as conjectural in the source reports need no
numerical evidence.  In their natural precise forms they are elementary
theorems.

\subsection{Total cumulants identify the geometric uniform law}

Let \(U_0,U_1,\ldots\) be independent random variables, each uniform on
\([-1,1]\), and put
\[
 X_{\mu,c,q}=\mu+c\sum_{j=0}^{\infty}q^jU_j,
 \qquad
 \mu\in\RealNumbers,\quad c>0,\quad 0<q<1.
\]
The series converges absolutely almost surely: with probability one all the
inequalities \(\lvert U_j\rvert\le1\) hold, and then its absolute value is
bounded by \(\sum_jq^j\).

\begin{theorem}[Total-cumulant identifiability]
\label{thm:total-cumulant-identifiability}
Let \(M\), \(V\), and \(K\) denote respectively the mean, variance, and
standardized fourth cumulant of \(X_{\mu,c,q}\), so
\(K=\kappa_4(X_{\mu,c,q})/V^2\).  Then
\[
 M=\mu,\qquad
 V=\frac{c^2}{3(1-q^2)},\qquad
 K=-\frac65\,\frac{1-q^2}{1+q^2}.
\]
The resulting map
\[
 (\mu,c,q)\longmapsto(M,V,K)
\]
is a bijection from
\(\RealNumbers\times(0,\infty)\times(0,1)\) onto
\(\RealNumbers\times(0,\infty)\times(-6/5,0)\).  Its inverse is
\[
 \mu=M,\qquad
 q=\sqrt{\frac{1+5K/6}{1-5K/6}},\qquad
 c=\sqrt{3V(1-q^2)}.
\]
\end{theorem}

\begin{proof}
For a uniform variable \(U\) on \([-1,1]\),
\[
 \Expectation U=0,\qquad
 \Expectation U^2=\frac13,\qquad
 \Expectation U^4=\frac15.
\]
Consequently its second and fourth cumulants are
\[
 \kappa_2(U)=\frac13,\qquad
 \kappa_4(U)=\Expectation U^4-3(\Expectation U^2)^2=-\frac{2}{15}.
\]
The partial sums converge to \(X_{\mu,c,q}\) in every \(L^r\), because their
tails are bounded deterministically by a geometric series.  Their moments
through order four therefore converge.  Independence, additivity of
cumulants, and homogeneity of \(\kappa_r\) give
\[
 \kappa_2(X_{\mu,c,q})
  =\frac{c^2}{3}\sum_{j\ge0}q^{2j}
  =\frac{c^2}{3(1-q^2)}
\]
and
\[
 \kappa_4(X_{\mu,c,q})
  =-\frac{2c^4}{15}\sum_{j\ge0}q^{4j}
  =-\frac{2c^4}{15(1-q^4)}.
\]
Dividing the latter by the square of the former yields the displayed formula
for \(K\).  Solving it gives
\[
 q^2=\frac{1+5K/6}{1-5K/6}.
\]
As \(K\) runs from \(-6/5\) to \(0\), this expression runs strictly from
\(0\) to \(1\).  The formulas for \(\mu\) and \(c\) then follow uniquely from
\(M\) and \(V\), and they plainly reconstruct the prescribed triple.
\end{proof}

Thus total cumulants identify the exact three-parameter model.  Stability
under noisy cumulants, missing observations, or a perturbed digit law is a
different problem; it becomes a theorem or conjecture only after the
perturbation model and metric are specified.

\subsection{A transcendental value determines the Gaussian rectangle}

For integers \(a,b\ge0\), write
\[
 G_{a,b}(Q)=\genfrac{[}{]}{0pt}{}{a+b}{a}_{Q}.
\]
This is the generating polynomial for partitions whose Ferrers diagrams fit
inside an \(a\)-by-\(b\) rectangle.

\begin{theorem}[Uniqueness at a transcendental base]
\label{thm:transcendental-gaussian-uniqueness}
Let \(q\in\ComplexNumbers\) be transcendental over \(\RationalNumbers\).  If
\[
 0<k<n,\qquad 0<k'<n',\qquad
 \genfrac{[}{]}{0pt}{}{n}{k}_{q}
   =\genfrac{[}{]}{0pt}{}{n'}{k'}_{q},
\]
then \(n'=n\) and either \(k'=k\) or \(k'=n-k\).
\end{theorem}

\begin{proof}
Transcendence says that evaluation at \(q\) is injective on
\(\IntegerNumbers[Q]\).  The asserted equality of values is therefore an equality of
Gaussian polynomials.  Put
\[
 a=k,\quad b=n-k,\quad
 m=\min(a,b),\quad M=\max(a,b),
\]
and define \(a',b',m',M'\) similarly.

The coefficient of \(Q^d\) in \(G_{a,b}\) counts partitions of \(d\) fitting
inside the rectangle.  Every partition of \(d\le m\) has at most \(m\) parts
and largest part at most \(m\), so these coefficients equal the unrestricted
partition numbers \(p(d)\).  At \(d=m+1\), at least one partition is excluded:
if \(a=m\), the partition \(1+\cdots+1\) with \(m+1\) parts is too tall; if
\(b=m\), the one-part partition \((m+1)\) is too wide.  Hence \(m+1\) is the
first degree at which the coefficient is smaller than \(p(d)\).  Equality of
the two polynomials therefore gives \(m=m'\).

The degree of \(G_{a,b}\) is \(ab=mM\), so equality of degrees and positivity
of \(m\) give \(M=M'\).  Thus the unordered pairs
\(\{a,b\}\) and \(\{a',b'\}\) agree.  Since \(n=a+b\) and \(k=a\), this says
exactly that \(n'=n\) and \(k'=k\) or \(k'=n-k\).
\end{proof}

\section{Precisely stated open frontiers}
\label{sec:precise-open-frontiers}

The following list keeps only assertions with a definite truth value.  Each
conjecture states its domain, normalization, and branch, followed by the
present evidence and the missing mathematical obligation.

\subsection{\texorpdfstring{Generic monodromy of finite \(q\)-products}{Generic monodromy of finite q-products}}

Set
\[
 P_{n,q}(a)=\prod_{j=0}^{n-1}(1-aq^j).
\]
The group-theoretic implication needed for this frontier is already a
theorem.

\begin{theorem}[Simple, distinct critical values force full monodromy]
\label{thm:simple-critical-values-monodromy}
Let \(P\in\ComplexNumbers[a]\) have degree \(n\ge2\).  Suppose \(P'\) has \(n-1\)
simple zeros and their images under \(P\) are pairwise distinct.  Then the
monodromy group of \(P(a)-y=0\) over
\(\ComplexNumbers\setminus\{\text{critical values of }P\}\) is \(S_n\).
\end{theorem}

\begin{proof}
Near a simple critical point, the holomorphic Morse lemma identifies the map
with \(z\mapsto z^2\) in local coordinates.  A small loop around its critical
value therefore exchanges precisely the two colliding sheets and fixes the
others: its monodromy is a transposition.  Distinct critical values give one
such branch cycle for each critical point.  Small positively oriented loops
around the finite critical values generate the fundamental group of their
complement; the loop at infinity is the inverse of their product.  Hence
these transpositions generate the monodromy group.

The cover is connected.  Algebraically, this follows because
\(P(a)-y\) is irreducible in \(\ComplexNumbers(y)[a]\): viewed in \(\ComplexNumbers[a,y]\), it has
degree one in \(y\), and any factor independent of \(y\) would divide the
coefficient \(-1\); Gauss's lemma then gives irreducibility over \(\ComplexNumbers(y)\).
Hence the monodromy action is transitive.

Associate a graph to the sheets, with an edge for every generating
transposition.  Its connected components are exactly the orbits of the group
they generate, so transitivity makes the graph connected.  The edge
transpositions of a connected graph generate all of \(S_n\).  Thus the
monodromy group is \(S_n\).
\end{proof}

\begin{conjecture}[Generic finite \(q\)-Pochhammer monodromy]
\label{conj:generic-pochhammer-monodromy}
For every \(n\ge3\) there is a nonzero polynomial \(E_n\in\IntegerNumbers[q]\) such that,
whenever \(q\in\ComplexNumbers\) and \(qE_n(q)\ne0\), the critical points of \(P_{n,q}\)
are simple and their critical values are pairwise distinct.  In particular,
the monodromy group is then \(S_n\).
\end{conjecture}

Simplicity is controlled by the discriminant of \(P'_{n,q}\), and coincidence
of two critical values by a second resultant or discriminant, all with integer
coefficients in \(q\).  The real obligation is to show that their product is
not the zero polynomial for every \(n\).  Once that is known, every excluded
nonzero parameter is algebraic.  A single exact specialization with simple
critical points and distinct critical values for each \(n\), or a uniform
symbolic argument, would suffice.  Floating-point plots cannot establish this
nonvanishing.

\subsection{Continuous Gaussian monotonicity}

For \(0<q<1\) and \(u>0\), use the standard normalization
\[
 \Gamma_q(u)
  =(1-q)^{1-u}
    \prod_{j=0}^{\infty}\frac{1-q^{j+1}}{1-q^{u+j}}.
\]
For \(x>y>0\), define the continuous Gaussian coefficient by
\[
 \mathcal G_q(x,y)=
 \begin{cases}
 \displaystyle
 \frac{\Gamma_q(x+1)}
      {\Gamma_q(y+1)\Gamma_q(x-y+1)},&0<q<1,\\[1.2ex]
 \displaystyle
 \frac{\Gamma(x+1)}
      {\Gamma(y+1)\Gamma(x-y+1)},&q=1,\\[1.2ex]
 \displaystyle
 q^{y(x-y)}\mathcal G_{1/q}(x,y),&q>1.
 \end{cases}
\]
The last line fixes the reciprocal-base normalization rather than leaving it
implicit.

\begin{conjecture}[Continuous Gaussian monotonicity]
\label{conj:continuous-gaussian-monotonicity}
For every pair \(x>y>0\), the function
\[
 q\longmapsto\mathcal G_q(x,y)
\]
is strictly increasing on \((0,1)\).
\end{conjecture}

For integral \(x=n\) and \(y=k\), this is immediate from the nonnegative
coefficients of the Gaussian polynomial, with strictness because
\(0<k<n\).  There is also rigorous local evidence at the upper endpoint.
Apply \cref{thm:qs-qgamma-one} at \(x+1\), \(y+1\), and \(x-y+1\).
Its linear coefficient satisfies
\[
 c_1(x+1)-c_1(y+1)-c_1(x-y+1)
 =-\frac{y(x-y)}2.
\]
Consequently, if \(q=\EulerE^{-t}\), then
\[
 \log\frac{\mathcal G_{\EulerE^{-t}}(x,y)}{\mathcal G_1(x,y)}
  =-\frac{y(x-y)}2t+O(t^2),
\]
so \(\mathcal G_{\EulerE^{-t}}(x,y)<\mathcal G_1(x,y)\) for every sufficiently
small \(t>0\).  A global proof must establish the sign of the
\(q\)-derivative of the displayed
\(q\)-gamma quotient, most likely by reorganizing the resulting Lambert
series into a manifestly positive expression.  Conversely, one
outward-rounded interval evaluation showing a negative derivative would be a
rigorous counterexample.

\subsection{A complete-Bernstein principal inverse}

Fix \(0<q<1\) and let
\[
 P_q(a)=(a;q)_\infty=\prod_{j=0}^{\infty}(1-aq^j),
 \qquad a<1.
\]
Every factor is positive.  Locally uniform differentiation of the logarithm
gives
\[
 \frac{P_q'(a)}{P_q(a)}
  =-\sum_{j=0}^{\infty}\frac{q^j}{1-aq^j}<0.
\]
Moreover, \(P_q(a)\to0\) as \(a\uparrow1\), while
\(P_q(a)\to\infty\) as \(a\to-\infty\).  For the latter limit, retain any
fixed number of factors and then let that finite product tend to infinity.
Thus \(P_q\) has a decreasing inverse
\[
 A_q:(0,\infty)\longrightarrow(-\infty,1),
 \qquad
 X_q(y)=1-A_q(y)>0.
\]
The first two Bernstein signs are theorems.  Indeed, if
\(r_j=q^j/(1-aq^j)>0\), then
\[
 \frac{P_q''(a)}{P_q(a)}
  =\left(\sum_{j\ge0}r_j\right)^2-\sum_{j\ge0}r_j^2
  =2\sum_{0\le i<j}r_ir_j>0.
\]
Implicit differentiation now gives
\[
 A_q'(y)=\frac1{P_q'(A_q(y))}<0,\qquad
 A_q''(y)
  =-\frac{P_q''(A_q(y))}{P_q'(A_q(y))^3}>0,
\]
and therefore \(X_q'(y)>0\) and \(X_q''(y)<0\).

\begin{conjecture}[Complete-Bernstein principal inverse]
\label{conj:complete-bernstein-principal-inverse}
For every \(0<q<1\), \(X_q\) is a complete Bernstein function.  Equivalently,
there exist \(\alpha_q,\beta_q\ge0\) and a positive measure \(\nu_q\) on
\((0,\infty)\), with
\[
 \int_{(0,\infty)}\frac{\nu_q(ds)}{1+s}<\infty,
\]
such that
\[
 X_q(y)=\alpha_q+\beta_qy+
        \int_{(0,\infty)}\frac{y}{y+s}\,\nu_q(ds)
 \qquad (y>0).
\]
\end{conjecture}

The conjecture implies all alternating derivative inequalities
\((-1)^{r-1}X_q^{(r)}(y)\ge0\) for \(r\ge1\).  The two proved signs above and
the higher signs observed in formal series reversion are evidence, not a
proof of the infinite family.  The decisive obligation is a Pick continuation
of \(X_q\) to the slit plane, or an explicit positive representing measure.
Notice also that the endpoint above forces
\(\alpha_q=X_q(0+)=0\) if the conjecture holds.

\subsection{\texorpdfstring{The minimum of the \(q\)-gamma function}{The minimum of the q-gamma function}}

The existence and uniqueness of the minimizer are not conjectural.

\begin{theorem}[Unique minimum of \(\Gamma_q\)]
\label{thm:qgamma-unique-minimum}
For every \(0<q<1\), the function
\(\Gamma_q:(0,\infty)\to(0,\infty)\) has a unique global minimizer
\(x_q^\ast\).
\end{theorem}

\begin{proof}
The logarithmic derivative of the standard \(q\)-gamma product is
\[
 \psi_q(x)
 =-\log(1-q)+(\log q)
   \sum_{j=0}^{\infty}\frac{q^{x+j}}{1-q^{x+j}}.
\]
The series and its derivative converge locally uniformly, and
\[
 \psi_q'(x)
  =(\log q)^2
    \sum_{j=0}^{\infty}
      \frac{q^{x+j}}{(1-q^{x+j})^2}>0.
\]
Thus \(\psi_q\) is strictly increasing.  Its \(j=0\) summand shows that
\(\psi_q(x)\to-\infty\) as \(x\downarrow0\), whereas dominated geometric
convergence gives
\(\psi_q(x)\to-\log(1-q)>0\) as \(x\to\infty\).
It therefore has exactly one zero \(x_q^\ast\).  Since
\(\Gamma_q'(x)=\Gamma_q(x)\psi_q(x)\) and \(\Gamma_q(x)>0\),
the function decreases before this zero and increases after it.  The zero is
the unique global minimizer.
\end{proof}

Write \(m_q=\Gamma_q(x_q^\ast)\).

\begin{conjecture}[Monotonic motion of the \(q\)-gamma minimum]
\label{conj:qgamma-monotone-minimizer}
On \(0<q<1\), the minimizer \(x_q^\ast\) is strictly increasing and the
minimum value \(m_q\) is strictly decreasing.
\end{conjecture}

The archived safeguarded-Newton computations, which are numerical evidence
rather than interval certificates, give
\[
\begin{array}{c|cccc}
q&0.05&0.20&0.60&0.95\\ \hline
x_q^\ast&1.37998346149&1.42196580342&1.45066979499&1.46060675866\\
m_q&0.98270463141&0.95376342969&0.91169489999&0.88840573828
\end{array}
\]
and exhibit both proposed signs.  They are compatible with the formal
\(q\downarrow0\) asymptotics and the classical-gamma limit as
\(q\uparrow1\).  A proof requires sign control of the implicit derivative
\[
 \frac{d x_q^\ast}{dq}
   =-\frac{\partial_q\psi_q(x_q^\ast)}
           {\partial_x\psi_q(x_q^\ast)}
\]
and, by the envelope identity,
\[
 \frac{d}{dq}\log m_q
   =\partial_q\log\Gamma_q(x_q^\ast).
\]
The denominators are already positive by the theorem; the two numerators
remain the substantive Lambert-series inequalities.

\subsection{Cyclotomic resurgence with fixed normalization}

Let \(\zeta\) be a primitive \(m\)-th root of unity and, for
\(\Re t>0\) and \(\lvert a\rvert<1\), take the logarithm determined by the
absolutely convergent series
\[
 L_a(t)
  =\PrincipalLogarithm(a;\zeta \EulerE^{-t})_\infty
  =-\sum_{r=1}^{\infty}
     \frac{a^r}{r(1-\zeta^re^{-rt})}.
\]
The singular action and constant term forced by this normalization are
\[
 S_{-1}(a)=-\frac{\operatorname{Li}_2(a^m)}{m^2},
\]
\[
 S_0(a)
  =\frac{1}{2m}\log(1-a^m)
   -\sum_{\substack{r\ge1\\m\nmid r}}
      \frac{a^r}{r(1-\zeta^r)}.
\]
A formal series \(\FormalGevreySeriesSymbol{R}(t)=\sum_{n\ge0}c_nt^n\) is
\emph{Gevrey one} here if
\(\lvert c_n\rvert\le CA^n n!\) for some \(C,A>0\).  Its formal Borel
transform is
\[
\mathscr B\FormalGevreySeriesSymbol{R}(\xi)
  =\sum_{n\ge0}\frac{c_{n+1}}{n!}\xi^n;
\]
the constant \(c_0\) is carried separately by Borel--Laplace summation.

\begin{conjecture}[Cyclotomic resurgent inverse]
\label{conj:cyclotomic-resurgent-inverse}
Let \(K\Subset\{a:\lvert a\rvert<1\}\) and
\(0<\delta<\pi/2\).  Uniformly for \(a\in K\), as \(t\to0\) in
\(\lvert\arg t\rvert\le\pi/2-\delta\), the normalized function
\[
 R_a(t)=L_a(t)-\frac{S_{-1}(a)}{t}-S_0(a)
\]
has a Gevrey-one asymptotic expansion.  Its Borel transform admits endless
analytic continuation: for each \(a\in K\), there is a locally finite set
\(\Sigma_{a,\zeta}\subset\ComplexNumbers\), with \(0\notin\Sigma_{a,\zeta}\), such that the
germ at zero continues along every path in
\(\ComplexNumbers\setminus\Sigma_{a,\zeta}\) issuing from zero.  In every direction whose
open Borel ray avoids \(\Sigma_{a,\zeta}\) and for which the Laplace integral
converges, its Borel sum equals the exact \(R_a(t)\).

Moreover, fix \(a_0\in\operatorname{int}K\), and let \(v(t)\to0\) be a
resurgent germ on a non-Stokes sector.  Let \(a(t)\in K\), with
\(a(t)\to a_0\), be the chosen local solution of
\[
 L_{a(t)}(t)=L_{a_0}(t)+v(t).
\]
If \(\partial_aL_{a(t)}(t)\) is bounded away from zero on every closed
subsector under consideration, then \(a(t)\) is resurgent, and its Stokes
automorphisms are obtained from those of \(L\) and \(v\) by the same local
implicit inversion.  Exponentiating the displayed equation gives the
corresponding statement for the chosen branch of the product inverse.
\end{conjecture}

The exact Lambert-series identity for \(L_a\), its residue-class splitting,
and the Bernoulli expansion of the \(m\)-divisible terms explain
\(S_{-1}\), \(S_0\), and the observed Gevrey growth.  They do not prove
endless continuation or locate every Borel singularity.  The remaining
obligations are uniform factorial remainder bounds on \(K\), analytic
continuation and growth estimates in the Borel plane, identification of the
Stokes set, equality of Borel sums with the convergent product, and a
quantitative implicit-function argument for the inverse clause.

\subsection{Editorial disposition of the remaining proposals}

Several former conjectures were deliberately not carried forward.

\begin{itemize}
\item The proposed finite-CDF inversion statement combined an elementary
  simple-fold theorem with an unproved analyticity assertion about a specific
  density.  The fold theorem is covered by the branch-aware certificates
  above; any density claim must name the density and regularity hypothesis.
\item The phrase \emph{generic Gaussian monodromy has the permitted group}
  did not define either the parameter space or the permitted group.  The
  finite-product conjecture above is its falsifiable replacement.
\item Broad inverse-universality and positivity-after-normalization claims did
  not specify a topology, scaling, normalization, or limiting class.
\item Total cumulant identifiability and discrete Gaussian uniqueness at a
  transcendental base were upgraded to
  Theorems~\ref{thm:total-cumulant-identifiability}
  and~\ref{thm:transcendental-gaussian-uniqueness}.
\end{itemize}

The omitted themes may still motivate research programs, but they are not
mathematical assertions until their quantifiers and normalizations are fixed.

\section{Lean source-status map}
\label{sec:lean-certification-map}

This section is a source-presence crosswalk, not compilation evidence.  The
declarations below occur in the indicated Lean source modules; no fresh build
status is asserted by this table.  A source declaration supports exactly the
topic stated in the last column; it does not silently formalize the stronger
complex certification theorems or conjectures of this chapter.  Module names
suppress the common
\nolinkurl{FabiusFunction.} prefix, and declaration names suppress the common
\nolinkurl{Fabius.} namespace; the two \nolinkurl{Bell.} declarations retain
their different namespace explicitly.  Every displayed suffix otherwise
matches the source exactly.

\begingroup
\small
\setlength{\tabcolsep}{4pt}
\renewcommand{\arraystretch}{1.18}
\begin{longtable}{@{}>{\raggedright\arraybackslash}p{0.21\textwidth}>{\raggedright\arraybackslash}p{0.45\textwidth}>{\raggedright\arraybackslash}p{0.29\textwidth}@{}}
\toprule
Lean module & Exact source declarations & Legitimate human crosswalk \\
\midrule
\endfirsthead
\toprule
Lean module & Exact source declarations & Legitimate human crosswalk \\
\midrule
\endhead
\midrule
\multicolumn{3}{r@{}}{\small\itshape Continued on the next page}\\
\endfoot
\bottomrule
\endlastfoot
\nolinkurl{FabiusInverse}
&
\nolinkurl{fabiusReal_fabiusInv};
\nolinkurl{fabiusInv_fabiusReal};
\nolinkurl{strictMonoOn_fabiusInv};
\nolinkurl{continuous_fabiusInv};
\nolinkurl{fabiusInv_hasDerivAt}
&
Two-sided inversion on the proved domain, strict monotonicity, continuity,
and the regular inverse derivative. \\

\nolinkurl{FabiusInverseEffectiveContinuity}
&
\nolinkurl{abs_fabiusInv_sub_lt_inverse_two_pow_of_lt_factorialDenominator};
\nolinkurl{fabiusInv_effectivelyUniformContinuous}
&
An explicit factorial-denominator inverse modulus and effective uniform
continuity. \\

\nolinkurl{FabiusInverseLogarithmicModulus}
&
\nolinkurl{abs_fabiusInv_sub_lt_inv_nat_of_lt_logarithmicFactorialDenominator};
\nolinkurl{fabiusInv_effectivelyUniformContinuous_logarithmic};
\nolinkurl{fabiusInv_effectivelyUniformContinuous_logarithmicDelta}
&
The later logarithmic inverse moduli, including the delta formulation. \\

\nolinkurl{LambertWBranchPointGeometry};
\nolinkurl{LambertWBranchPointAsymptotics}
&
\nolinkurl{tendsto_deriv_principalLambertW_branchPoint_atTop};
\nolinkurl{tendsto_deriv_lowerLambertW_branchPoint_atBot};
\nolinkurl{principalLambertW_not_differentiableAt_branchPoint};
\nolinkurl{lowerLambertW_not_differentiableAt_branchPoint};
\nolinkurl{principalLambertW_add_one_isEquivalent_branchPoint};
\nolinkurl{lowerLambertW_add_one_isEquivalent_branchPoint}
&
The opposite vertical tangents, failure of finite differentiability, and the
signed leading square-root laws at the real Lambert-\(W\) branch point.  They
identify the correct signed leading square-root scale for a possible
regularization; they establish neither a regularized analytic coordinate
chart nor a higher Puiseux remainder bound. \\

\nolinkurl{GeometricUniformLaw}
&
\nolinkurl{hasSum_geometricUniformWeight};
\nolinkurl{geometricUniformSeries_mem_Icc};
\nolinkurl{geometricUniformDistribution_selfSimilar};
\nolinkurl{geometricUniformDistribution_support_eq_Icc}
&
Convergence, support, and self-similarity of Lean's normalized \([0,1]\)
geometric uniform law.  If
\(Z=(1-q)\sum_{j\ge0}q^jV_j\) and \(U_j=2V_j-1\), then
\(\sum_{j\ge0}q^jU_j=(2Z-1)/(1-q)\).  This affine rescaling underlies
Theorem~\ref{thm:total-cumulant-identifiability}. \\

\nolinkurl{GeometricSincCharacteristicFunction}
&
\nolinkurl{charFun_geometricUniformDistribution_eq_phase_mul_geometricSincProduct};
\nolinkurl{charFun_geometricUniformDistribution_eq_phase_mul_geometricReciprocalGamma}
&
The sinc-product characteristic function and its reciprocal-gamma form. \\

\nolinkurl{SincZetaRemainder};
\nolinkurl{FabiusComputability}
&
\nolinkurl{norm_sincZeta_tail_le};
\nolinkurl{rvachevFourierProduct_eq_prefix_mul_correction_mul_tail};
\nolinkurl{abs_uniformPartialCDF_sub_fabiusReal_le};
\nolinkurl{fabiusEffectiveUniformModulus}
&
Certified sinc-product and CDF tails and an effective forward Fabius modulus.
These are nearby formal certificates, not the general-purpose complex
Theorems~\ref{thm:qpoch-tail-certificate}--\ref{thm:validated-argument-principle}. \\

\nolinkurl{QPochhammerElementaryIdentities};
\nolinkurl{QBinomialReciprocity}
&
\nolinkurl{finiteQPochhammerIn_base_reversal};
\nolinkurl{finiteQPochhammerIn_inv_base_reversal};
\nolinkurl{gaussianBinomial_adjacent_mul};
\nolinkurl{gaussianBinomial_reciprocity};
\nolinkurl{gaussianBinomial_neg_one_eq_zero_of_odd_degree};
\nolinkurl{gaussianBinomial_neg_one_even_odd_eq_zero}
&
Finite base reversal, adjacent ratios, Gaussian reciprocity, and the
odd-degree vanishing mechanism at \(q=-1\). \\

\nolinkurl{GaussianBinomialAtNegOne}
&
\nolinkurl{gaussianBinomial_neg_one_even_even};
\nolinkurl{gaussianBinomial_neg_one_odd_even};
\nolinkurl{gaussianBinomial_neg_one_odd_odd};
\nolinkurl{finiteQPochhammerIn_neg_one_even};
\nolinkurl{finiteQPochhammerIn_neg_one_odd}
&
The remaining parity evaluations for Gaussian coefficients and the finite
\(q\)-Pochhammer formulas at \(q=-1\); together with the preceding row these
form the complete parity atlas. \\

\nolinkurl{GaussianBinomialAtNegOneDerivative}
&
\nolinkurl{gaussianBinomial_derivative_eval_neg_one_of_even_degree};
\nolinkurl{gaussianBinomial_derivative_eval_neg_one_even_odd};
\nolinkurl{gaussianBinomial_even_odd_rootMultiplicity_int};
\nolinkurl{gaussianBinomial_even_odd_rootMultiplicity}
&
The complete first-jet layer at \(q=-1\): two commutative-ring derivative
formulas, plus multiplicity one for the alternating even-row/odd-column zero
over \(\IntegerNumbers\) and every characteristic-zero commutative ring when
the column is admissible. \\

\nolinkurl{BellPolynomialInversion};
\nolinkurl{QuadraticAsymptoticInversion}
&
\nolinkurl{Bell.complete_cumulant};
\nolinkurl{Bell.cumulant_complete};
\nolinkurl{quadratic_asymptotic_inversion}
&
Mutual moment--cumulant inversion and the proved quadratic asymptotic
inversion schema. \\
\end{longtable}
\endgroup

No declaration in this audited source crosswalk is cited as proving the exact
statements of the general
\(q\)-Pochhammer derivative-tail theorem, interval Newton theorem, complex
Krawczyk theorem, Rouch\'e disk theorem, validated argument principle,
total-cumulant parameter-identification theorem, or transcendental Gaussian
uniqueness theorem as stated here.  They are complete human-readable proofs
in this chapter, not claimed Lean milestones.  Likewise, no Lean counterpart
is claimed for the five conjectures in
Section~\ref{sec:precise-open-frontiers}.
